\begin{abstract}%
Physics-informed neural networks (PINNs) can be accurate on average yet inaccurate
in thin layers and at moving interfaces. Useful uncertainty estimates should
reflect these regional differences. We ask whether post-hoc Gaussian uncertainty
from Gauss--Newton curvature, a linearised Laplace approach, places uncertainty
where numerical errors occur. Its covariance depends on a prior scale and a
common residual-noise scale. We evaluate full-network covariances for $17$
Burgers and Allen--Cahn runs with two
architectures, comparing reference-defined sharp and bulk regions. Allocation is
strongly scale-dependent. For the trained prediction, five of eight Burgers--MLP
runs have identified precisions matching sharp-to-bulk uncertainty and error
ratios; these span more than four orders of magnitude. No match is found for three
runs in the tested range. Neither training-based scale choices nor the tested
predictive centres reliably calibrate both regions. Even at the matching points,
a reference-informed oracle calibrating bulk coverage to $68\%$ yields only
$30.3$--$55.7\%$ sharp coverage for the trained prediction. Reference-based readout
refits reveal features capable of representing sharp structures. Spectral analysis
explains how the scales redistribute uncertainty across curvature modes, with key
bounds machine-checked in Lean~4. Matching regional uncertainty ratios therefore
does not ensure calibrated intervals: allocation and coverage must be assessed
separately.
\end{abstract}

\begin{keywords}
physics-informed neural networks, uncertainty quantification, Gauss--Newton covariance,
calibration, spectral analysis
\end{keywords}
